En una experiència d'interferències sonores es fan servir dos altaveus, separats 6~m l'un de l'altre, que generen una ona acústica coherent i monocromàtica de freqüència 170~Hz. Observa que un altaveu esta situa a $(-3, 0)$ i l'altre a $(3,0)$.

\figura[0.45]{fig1.png}

\begin{apartat}{1,25}
Determina la longitud de l'ona generada. Indica, pels punts representats a la figura, en quins punts la intensitat sonora serà màxima i en quins punts mínima. Justifica al resposta.
\end{apartat}

\begin{apartat}{1,25}
Per determinar amb claredat en quins punts la interferència és màxima, el professor demana als estudiants que es tapin una orella. Per quina raó cal fer-ho? Que passarà amb la intensitat sonora al punt $(0,2)$ si introduïm un petit retard temporal entre el senyal emès per un dels altaveus respecte a l'altre? Justifica les respostes.
\end{apartat}

\dades{La velocitat del so a l'aire és de 340~m/s.}
